\subsection{Forward rates, expected returns and curve trades}

The examples in \hyperref[meth:forward-rates-expected-returns]{forward rates and bond returns} use annual compounding unless stated, and the curves they use are built in \hyperref[meth:spot-par-forward-curves]{spot, par and forward rates}. Their numbers are checked in \texttt{notes/artefacts/term-structure-dynamics/term-structure-checks.py}.

\begin{qfmethod}[Forward rates and bond returns: forwards as break-evens, the rolling yield, realised against expected returns, assumptions about how spot rates evolve relative to forwards, and riding the yield curve]\label{meth:forward-rates-expected-returns}
\qffield{Statement} \emph{Break-evens.} The implied spot rates one year forward are the future spot rates at which every bond earns the same one-year return as the one-year zero. If forwards are realised, all government bonds earn the same holding-period return and every self-financed curve position breaks even. If the curve is unchanged, an $n$-year zero earns the forward $f_{n-1,n}$, its \emph{rolling yield}: yield income plus the \emph{rolldown} gain as it ages down an upward-sloping curve \citep[chap.~31, pp.~695 and~711--712]{s7-fabozzifixed}. Approximately, $f_{n-1,n}\approx s_n+(n-1)(s_n-s_{n-1})$ (eq.~(31C.2)). Expected return decomposes into yield income, rolldown, the value of convexity and the duration effect of an expected curve change (chap.~32, pp.~715 and~724--725).

\emph{Assumptions in active management.} A manager who expects forwards to be realised expects no bond to beat cash. One who expects an unchanged curve (rates follow a random walk) expects longer bonds to earn their forwards, and so to beat cash when the curve slopes up. Ilmanen reports that the unchanged-curve forecast has been more accurate than the forwards on average, and that steeper curves have been followed by larger excess returns, against the pure expectations hypothesis (pp.~702--703 and~712).

\qffield{Practice} \emph{Reproduced.} Ilmanen's example: one- and two-year spot rates of 6 and 8.08 percent imply $f_{1,2}=10.20$ percent; if it is realised, the two-year zero rises from 85.61 to 90.74, a 6 percent return, the same as the one-year zero. The one-year rate must rise by more than 420 basis points before the two-year zero underperforms (p.~695). His approximation gives 10.16, 13.00, 14.28 and 13.76 percent against exact forwards of 10.20, 13.07, 14.36 and 13.77 (p.~713).

\emph{Original} (the curve of \hyperref[meth:spot-par-forward-curves]{spot, par and forward rates}). If forwards are realised, the two- to five-year zeros all earn 3.80 percent in the first year; if the curve is unchanged, each $n$-year zero earns $f_{n-1,n}$. Riding the curve, buying the 5-year zero and selling it after a year, returns 4.5235 percent if the curve is unchanged against 3.80 percent on the 1-year zero, and 2.5423 percent if all rates rise 50 basis points. The ride breaks even at an 18.1 basis point rise in the 4-year rate, the forward-implied change $f_{1,5}-s_4$.

\qffield{Conditions} The identities are exact for zero-coupon bonds and default-free cash flows; for coupon bonds the rolling yield and forwards differ slightly because of reinvestment assumptions (chap.~32, p.~724, n.~7). Realised returns depend on the actual path of rates, reinvestment and the sale price (\hyperref[meth:bond-return-decomposition]{sources of return}). Expected returns need a model of term premia and convexity; neither is observable (chap.~31, pp.~701--702).

\qffield{Pitfalls} Taking a bearish position whenever rates are expected to rise: it pays only if they rise by more than the forwards imply, because the long bond's positive carry is the cushion (pp.~703--704). Treating rolldown as risk-free income: it assumes the curve does not change. Calling a forward an expectation of the future spot rate (\hyperref[res:term-structure-theories]{theories}).

\qffield{See also} \hyperref[meth:curve-trades-carry-rolldown]{curve trades, carry and rolldown}; \hyperref[def:bond-yield-spread-measures]{yield to maturity}.

\end{qfmethod}

\begin{qfmethod}[Curve trades, carry and rolldown, and inversion as a recession signal: DV01-neutral steepeners and flatteners, bull and bear moves, carry signs that depend on the curve, and the Estrella--Trubin record]\label{meth:curve-trades-carry-rolldown}
\qffield{Statement} Flatteners and steepeners trade the spread between two points of the curve, ``done DV01 neutral'', for example 2s10s, the 10-year yield minus the 2-year \citep[chap.~9, sec.\ ``Flatteners and Steepeners'', PDF pp.~266--267]{s7-gliner}; a \emph{bear flattener} has short rates rising more than long rates, a \emph{bull flattener} long rates falling more than short (p.~267, Figure~9.7), and by symmetry a \emph{bull steepener} has short rates falling more and a \emph{bear steepener} long rates rising more. A steepener is long the short maturity and short the long maturity in DV01-matched amounts.

\emph{Carry} is the income of a position net of financing over a holding period; \emph{rolldown} is the price change as a bond ages along an unchanged curve. Total return $\approx$ carry $+$ rolldown $+$ the effect of yield changes.

\emph{Inversion as a recession signal.} Estrella and Trubin's preferred spread is the 10-year constant-maturity Treasury yield minus the three-month secondary-market bill rate on a bond-equivalent basis, in monthly averages; it ``has turned negative before each recession in the period from January 1968 to July 2006''. All six NBER recessions since 1968 were preceded by at least three negative monthly averages in the twelve months before, with no false signals on a monthly-average basis, while negative spreads occurred on 100 days between 1968 and 2005 in months that did not have a negative average \citep[pp.~4--5]{estrella-trubin-2006}. The 2s10s spread ``tends to turn negative earlier and more frequently'' (p.~4).

\qffield{Notation} DV01 price change per 1~bp yield change (\hyperref[mod:duration-convexity]{duration and convexity}); a DV01-neutral steepener holds $N_2\,\mathrm{DV01}_2=N_{10}\,\mathrm{DV01}_{10}$.

\qffield{Conditions} \emph{Carry signs.} carry ``$+$ for steepeners, $-$ for flatteners in normal yield curve'' is not general. For a DV01-neutral steepener (long 2-year, short 10-year), carry per unit DV01 is positive with 2-year at 1.5 percent, 10-year at 3 percent and repo at 1 percent, but negative with an inverted front (2-year 4.5, 10-year 4.0, repo 5.0) and also negative with a mildly upward curve when repo is close to the 2-year yield (2.0, 2.5, 2.4). Rolldown is deterministic only if the curve is unchanged, and carry is known at inception only for a fixed financing rate.

\qffield{Practice} Gliner defines rolldown as rolling a futures or bond position to keep the original maturity profile (p.~268); that is the roll of a position, not rolldown in the usual sense defined above. Porter's view that a steepener is in effect a bet on a bond rally (\hyperref[meth:macro-practitioner-frameworks]{practitioner frameworks}) reflects the empirical tendency of the curve to steepen in rallies. Gliner reads an inversion far out the curve, as in Portugal and Greece in 2010--2012, as pricing default risk, and a mild inversion near the front as pricing cuts (sec.\ ``Curve Inversion'', p.~271; \hyperref[def:sovereign-credit-defaults]{sovereign credit}).

\qffield{Pitfalls} \emph{Three different statements}. (1) \emph{A historical observation}: for each of the six NBER recessions from 1968 to 2005, the monthly-average spread was negative in at least three of the twelve months before the recession began (3 to 12 such months; minima from $-0.08$ to $-3.51$) (table, p.~6). (2) \emph{The paper's forecast}: a probit for recession \emph{status} twelve months ahead, $P(\text{recession in month }t+12)=F(-0.6045-0.7374\,s_t)$, estimated on 1959--2005 data, where recession months run from the month after an NBER peak through the trough (pp.~2--3); a zero spread gives 27 percent. (3) \emph{Onset rules}, defined here and not in the paper, applied to negative months during expansions: a \emph{month rule} (a recession begins in $t+1$ to $t+12$) and an \emph{episode rule} (a run of consecutive negative months is followed by a recession beginning within twelve months of its last month).

\emph{The episode rule is retrospective}. It is scored from an episode's \emph{last} inverted month, which is known only after the spread has turned positive again, and against NBER peak dates announced months to more than a year after the event; at a date inside an episode neither the end of the run nor the recession's start is observable. Its hit rates below therefore describe the historical record and are not a prospective forecasting result; nor is the month rule one, since it too is scored on ex-post NBER dates, though each of its months is observable when it occurs. Neither is comparable with the probit's probability, which the paper states as a forecast made at $t$.

The paper's statement that monthly-average inversions gave ``no false signals'' (p.~5) matches the episode rule, not the month rule: its own footnote gives $-0.16$ in June 1989, ``fourteen months before the recession'' of August 1990 (p.~6, n.~a to the table). Replicating the paper's series from the Federal Reserve's H.15 data (three-month secondary-market bill converted to a bond-equivalent yield by the paper's formula, n.~4, against the 10-year CMT; business-day averages by month) \citep[H.15 data package, series RIFSGFSM03\_N.B and RIFLGFCY10\_N.B]{fed-h15-data} reproduces the table's counts and minima (the 1990 minimum $-0.07$ against the printed $-0.08$) and the June 1989 value; over 1968--2005 the 54 negative expansion months include 6 month-rule misses (December 1968, November 1978 to January 1979, June and July 1989), while all 9 episodes satisfy the episode rule.

So (1) and the episode-level claim hold in sample; no general twelve-month onset law follows from six recessions. A typical lead time of 6--18 months matches none of the three rules. \emph{After the paper's sample}, on the paper's series: the 43 complete negative expansion months from 2006 (October 2026, with data to 14 October, flagged incomplete and excluded) form five episodes, August 2006--May 2007, May--September 2019, February 2020, November 2022--November 2024 and March--April 2025. NBER peaks after 2005 are December 2007 and February 2020 (first contraction months January 2008 and March 2020), and no later peak was dated as of 16 October 2026 \citep[chronology table and dating page]{nber-business-cycles}.

The month rule gives 11 hits, 5 misses (August--December 2006, 13 to 17 months before January 2008) and 27 provisional misses from November 2022; the episode rule gives three hits (the 2020 one before a pandemic-driven recession) and two provisional misses. \emph{Recession outcomes through 17 October 2026}. Two conditions matter. \emph{Calendar}: the twelve-month window of every one of the 27 months, the last ending in April 2026, has now passed.

\emph{Dating}: NBER still lists February 2020 as the latest peak (chronology data last updated 14 March 2023; latest announcement 19 July 2021, the April 2020 trough) \citep[chronology table and announcements page]{nber-business-cycles}. Its peak announcements have come 4 to 12 months after the peak: January 1980 on 3 June 1980, July 1981 on 6 January 1982, July 1990 on 25 April 1991, March 2001 on 26 November 2001, December 2007 on 1 December 2008, February 2020 on 8 June 2020. So the 24 months from November 2022 to October 2024 are misses even allowing the longest lag yet seen; a recession dated in their windows would need a longer one. November 2024, March 2025 and April 2025, and with them both post-2022 episodes, are misses only if no peak up to April 2026 is announced within the lags seen so far.

Retrospective dating cannot be excluded for any of them. These outcomes stay provisional, for the dating reason only. Treasury's 3-month CMT, a proxy and not the paper's series, gives the same picture with 46 months and a further episode in June--August 2025 \citep[daily CSV files]{ust-par-yield-curve}. Both monthly series are rebuilt from daily data, with unique dates and consecutive months.

\qffield{See also} \hyperref[mod:duration-convexity]{duration and convexity}; \hyperref[mod:forward-rates-fra]{forward rates and term-structure theories}; \hyperref[def:macro-regime-asset-map]{regime map}; \hyperref[conv:government-bond-futures]{government bond futures}.

\qffield{Further reading} Term-structure decomposition \citep{rt-gurkaynak-wright-2012}.

\end{qfmethod}
